%%%%%%%%%%  scan idx 8  |  folio VII  %%%%%%%%%%
\section*{\underline{Introduction}}
Le but du présent Séminaire est de développer une théorie \underline{globale} des intersections, et une formule de Riemann--Roch, pour des schémas quelconques. Le lecteur trouvera une description du programme du Séminaire dans l'exposé 0. La possibilité en principe d'une démonstration d'une formule de Riemann--Roch pour les schémas \underline{réguliers} généraux, suivant les lignes du rapport bien connu de Borel--Serre en 1958, était claire dès ce moment, du moins pour un morphisme projectif. L'extension au cas général est moins évidente~; le programme dans lequel elle s'insère (et partiellement réalisé dans notre séminaire) remonte à 1960. Comme c'était également le cas pour la théorie de dualité des faisceaux cohérents (cf.\ R.~Hartshorne, Residues and Duality, Lecture Notes in Mathematics n\textsuperscript{o}~20, Springer), un outil essentiel pour formuler une théorie satisfaisante est la théorie des \underline{catégories dérivées} de Verdier, dont la connaissance est indispensable pour la compréhension du Séminaire.

À part cette théorie, l'étude du Séminaire n'exige guère qu'une connaissance générale des fondements de la Géométrie Algébrique, tels qu'ils sont exposés dans EGA, chapitres I, II, III~; nous n'aurons en plus que des références occasionnelles à faire à EGA IV, pour certains faits concernant les morphismes plats ou lisses, qui pour l'essentiel figurent également dans les premiers exposés de SGA~1. Enfin, pour développer certains résultats avec toute la généralité souhaitable pour les applications, nous faisons usage parfois de la notion de \underline{site annelé} et de \underline{topos annelé}, pour laquelle nous renvoyons à SGA~4, exposés I à IV. Le lecteur qui ignorerait le langage des sites et topos pourra remplacer partout lesdits objets par des espaces topologiques ordinaires, les objets du topos étant alors remplacés par des ouverts de ces espaces~; mais nous lui conseillons néanmoins, de préférence, de s'assimiler le langage des topos, qui fournit un principe d'unification extrêment commode.

La notion fondamentale pour la théorie présentée ici est celle de \underline{complexe de Modules parfait}, et ses diverses variantes, développées dans les exposés I, II, III. Il semble clair que ces notions, importantes également dans d'autres contextes en Géométrie Algébrique (notamment pour les formules de type Lefschetz--Verdier en cohomologie à coéfficients discrets (SGA 5) ou cohérents), auront aussi leur rôle à jouer en dehors de la Géométrie Algèbrique, notamment pour la formulation d'une variante analytique com%
